\documentclass[12pt]{article}
\usepackage[utf8]{inputenc}
\usepackage[T1]{fontenc}
\usepackage{hyperref}
\title{Aore Language Resources for K-12 Curriculum Development}
\author{Compiled by Claude Sonnet 5.5 (claude-sonnet-5-5) for Andr\'as Kornai}
\date{}
\begin{document}
\maketitle

\section{Introduction}
Aore (ISO 639-3 aor) is an Oceanic language formerly spoken on Aore Island, just off Espiritu Santo, Vanuatu. It is classified as extinct (UNESCO) or dormant (Endangered Languages Project). A 1972 survey by Tryon reported a single elderly speaker; Ethnologue gives extinction in the 1980s.

\section{Primary Educational Resources}
None found.

\section{Linguistic and Reference Materials}
\begin{itemize}
\item Wikipedia: \url{https://en.wikipedia.org/wiki/Aore_language}
\item Glottolog (aore1237): \url{https://glottolog.org/resource/languoid/id/aore1237.bigmap.html}
\end{itemize}

\section{Religious and Community Resources}
None found.

\section{Digital and Audio Resources}
Endangered Languages Project entry: \url{https://endangeredlanguages.com/lang/3851}

\section{Curriculum Development Context}
The search found no grammar, dictionary or recordings of Aore. A heritage curriculum would have to rely on archival work and on the related living languages of South Santo, and on local community interest on Aore Island, none of which was documented in the results.
\end{document}